\section{Shell编辑}
\subsection{快捷键}
\begin{table}[H]
	\centering
	\begin{tabular}{|l|l|}
		\hline
		\textbf{快捷键} & \textbf{说明} \\
		\hline
		Ctrl + T     & 打开新标签页。     \\
		\hline
		Ctrl + W     & 关闭标签页。      \\
		\hline
		Ctrl + L     & 跳转地址栏。      \\
		\hline
		Ctrl + Tab   & 切换标签页。      \\
		\hline
	\end{tabular}
	\caption{浏览器 常用快捷键}
\end{table}

\begin{table}[H]
	\centering
	\begin{tabular}{|l|l|}
		\hline
		\textbf{快捷键}   & \textbf{说明}      \\
		\hline
		Ctrl + a       & 移动光标到行首。         \\
		\hline
		Ctrl + e       & 移动光标到行尾。         \\
		\hline
		Alt + f        & 将光标向前移动一个单词。     \\
		\hline
		Alt + b        & 将光标向后移动一个单词。     \\
		\hline
		Ctrl + k(kill) & 剪切从光标位置到行尾的内容。   \\
		\hline
		Ctrl + y(yank) & 粘贴最近删除的内容。       \\
		\hline
		Ctrl + l       & 清屏，相当于执行clear命令。 \\
		\hline
		Ctrl + r       & 反向搜索历史命令。        \\
		\hline
	\end{tabular}
	\caption{Bash 常用快捷键}
\end{table}

\subsection{通配符}
注意通配符是shell处理的，而不是命令处理的。
*并不会匹配以.开头的文件，要匹配这些文件需要使用.[!.]*。
\begin{table}[H]
	\centering
	\begin{tabular}{|l|p{7cm}|}
		\hline
		\textbf{通配符}           & \textbf{说明}                                        \\
		\hline
		\texttt{*}             & 匹配任意多个字符。                                          \\
		\hline
		\texttt{?}             & 匹配任意单个字符。                                          \\
		\hline
		\texttt{[characters]}  & 匹配方括号内的任意单个字符。                                     \\
		\hline
		\texttt{[!characters]} & 匹配不在方括号内的任意单个字符。                                   \\
		\hline
		\texttt{[[:class:]]}   & 匹配属于字符类的任意单个字符，字符类有alnum、alpha、digit、lower、upper等。 \\
		\hline
	\end{tabular}
	\caption{常见通配符与说明}
\end{table}

\subsection{常用命令}
\begin{table}[H]
	\centering
	\begin{tabular}{lllllll}
		alias    & bg     & bind & builtin & cd      & command & compgen \\
		complete & disown & echo & enable  & exec    & exit    & export  \\
		fc       & fg     & hash & help    & history & jobs    & kill    \\
		logout   & printf & pwd  & type    & ulimit  & umask   & wait    \\
	\end{tabular}
	\caption{Bash 内置命令}
\end{table}

\begin{itemize}
	\item 查找命令：apropos、dpkg -S。
	\item 帮助命令：help、info、man、whatis、whereis、which。
	\item 编辑命令：wc、head、tail、cut、grep、sort、uniq。
	\item 输出重定向：\verb|> file 2>&1|相当于\verb|&> file|，|管道将命令的标准输出传给下一个命令的标准输入。
	\item 输入重定向：cat读取文件到标准输出，tee读取标准输入并将其同时写入标准输出和文件中，xargs将标准输入转换为命令的参数。
\end{itemize}

\begin{code}[title=常用命令示例]
	\begin{verbatim}
date                              # 显示当前日期时间
wc -l file                        # 统计行数（-w词数 -c字节数）
head -n 5 file                    # 显示前5行
tail -n 5 file                    # 显示后5行
tail -f file                      # 实时跟踪文件末尾（适合看日志）
ls /bin /usr/bin | sort | uniq | less  # 管道：合并排序去重分页
export PATH=~/bin:"$PATH"             # 将 ~/bin 追加到 PATH
	\end{verbatim}
\end{code}

\subsection{shell 扩展}
参数扩展、算术扩展和命令替换仍然会在双引号中进行，但单引号中所有扩展均不能用。
\begin{itemize}
	\item 参数扩展：\verb|${var}|表示变量的值，\verb|${#var}|表示变量值的长度。
	\item 算术扩展：\verb|$(( expression ))|，如\verb|echo $(( 3 + 5 ))|。
	\item 命令替换：\verb|$( command )|，如\verb|echo "Today is $(date)"|。
	\item 进程替换：\verb|<(cmd)|可以将命令的输出当成文件使用。
	\item 路径名扩展：\verb|~/|表示用户主目录，\verb|~+|表示当前工作目录，\verb|~-|表示上一个工作目录。
	\item 花括号扩展：\verb|{1..5}|表示生成1到5的数字序列，\verb|{a..e}|表示生成a到e的字母序列。
	\item 历史扩展：\verb|!!|表示上一条命令，\verb|!n|表示第n条命令，\verb|!string|表示最近一条以string开头的命令，
	      \verb|!?string|表示最近一条包含string的命令。如果你需要记录shell对话，可以使用\verb|script|命令。
\end{itemize}

\subsection{正则表达式}
正则表达式中存在元字符（metacharacter）:\verb@. ? + * { } ^ $ [ ]  - ( ) | \@

POSIX 字符类可用于方括号表达式中，常用的有：\texttt{[:upper:]}（大写字母）、\texttt{[:lower:]}（小写字母）、\texttt{[:alpha:]}（字母）、\texttt{[:digit:]}（数字）、\texttt{[:alnum:]}（字母和数字）、\texttt{[:space:]}（空白字符）。例如 \verb|[[:upper:]]*| 匹配以大写字母开头的任意字符串。

\begin{code}[title=简写字符类（大写为补集）]
	\begin{verbatim}
\d  [[:digit:]]    # 数字           \D  非数字
\s  [[:space:]]    # 空白字符       \S  非空白
\l  [[:lower:]]    # 小写字母       \L  非小写
\u  [[:upper:]]    # 大写字母       \U  非大写
\w  [_[:alnum:]]   # 单词字符       \W  非单词字符
\b  单词边界                        \B  非单词边界
	\end{verbatim}
\end{code}

\begin{code}[title=锚点与量词]
	\begin{verbatim}
^       行首（在 [...] 中为取反）
$       行尾
^$      空行

?       0 或 1 次
*       0 或多次
+       1 或多次
{n}     恰好 n 次
{n,}    至少 n 次
{n,m}   n 到 m 次（含两端）
	\end{verbatim}
\end{code}

\subsection{历史记录}
ctrl+r 搜索历史记录，esc退出搜索。

\begin{code}[title=history]
	\begin{verbatim}
history 10     # 显示最近 10 条历史记录
history -c     # 清空历史记录
!-3:p          # 打印倒数第 3 条命令（不执行）
rm !$          # !$ 表示上一条命令的最后一个参数，先 ls 再 rm !$
^jg^jpg        # 将上一条命令中的 jg 替换为 jpg 并执行（只替换首次出现）
!!:s/jg/jpg/   # 同上，替换语法的另一种写法
	\end{verbatim}
\end{code}

\subsection{颜色}
Shell中也能显示颜色，分为文本颜色和背景颜色，关闭颜色的代码是\verb|\033[0m|。
\begin{table}[H]
	\centering
	\begin{minipage}{0.45\linewidth}
		\centering
		\begin{tabular}{|l|l|}
			\hline
			\textbf{转义代码}     & \textbf{文本颜色}    \\
			\hline
			\verb|\033[0;30m| & 黑色 (black)       \\
			\verb|\033[0;31m| & 红色 (red)         \\
			\verb|\033[0;32m| & 绿色 (green)       \\
			\verb|\033[0;33m| & 棕色 (brown)       \\
			\verb|\033[0;34m| & 蓝色 (blue)        \\
			\verb|\033[0;35m| & 紫色 (magenta)     \\
			\verb|\033[0;36m| & 青色 (cyan)        \\
			\verb|\033[0;37m| & 浅灰色 (light grey) \\
			\hline
		\end{tabular}
		\caption{常用文本颜色代码}
	\end{minipage}
	\hfill
	\begin{minipage}{0.45\linewidth}
		\centering
		\begin{tabular}{|l|l|}
			\hline
			\textbf{转义代码}     & \textbf{文本颜色}       \\
			\hline
			\verb|\033[1;30m| & 深灰色 (deep grey)     \\
			\verb|\033[1;31m| & 浅红色 (light red)     \\
			\verb|\033[1;32m| & 浅绿色 (light green)   \\
			\verb|\033[1;33m| & 黄色 (yellow)         \\
			\verb|\033[1;34m| & 浅蓝色 (light blue)    \\
			\verb|\033[1;35m| & 浅紫色 (light magenta) \\
			\verb|\033[1;36m| & 浅青色 (light cyan)    \\
			\verb|\033[1;37m| & 白色 (white)          \\
			\hline
		\end{tabular}
		\caption{高亮文本颜色代码}
	\end{minipage}
\end{table}

\begin{table}[H]
	\centering
	\begin{tabular}{|l|l|}
		\hline
		\textbf{转义代码}   & \textbf{背景颜色} \\
		\hline
		\verb|\033[40m| & 黑色            \\
		\verb|\033[41m| & 红色            \\
		\verb|\033[42m| & 绿色            \\
		\verb|\033[43m| & 棕色            \\
		\verb|\033[44m| & 蓝色            \\
		\verb|\033[45m| & 紫色            \\
		\verb|\033[46m| & 青色            \\
		\verb|\033[47m| & 浅灰色           \\
		\hline
	\end{tabular}
	\caption{背景颜色代码}
\end{table}
